% GENERATED FILE. Edit canonical lessons, then run npm run generate:books.
% canonical-source-hash: fnv1a64:3e203b01d71ab8ef
% canonical-lessons: ML-C437-ganitam, ML-C437-sampattikasastram, ML-C437-tatvasastram, ML-C437-vijnanam, ML-C437-pathanam

\chapter{Subjects}
\label{ch:ml-subjects}
\hlchaptermodality{437}

\begin{chapteropening}
\textbf{By the end of this chapter:} \emph{I can say mathematics, economics, philosophy, knowledge and a study.}

Complete the last lesson of chapter 437: I can say mathematics, economics, philosophy, knowledge and a study.
\end{chapteropening}

\section[\texorpdfstring{\ml{ഗണിതം} (gaṇitaṁ)}{ഗണിതം (gaṇitaṁ)}]{\ml{ഗണിതം} (gaṇitaṁ) — formal mathematics; maths}

\label{lesson:ML-C437-ganitam}



\begin{quote}
Before the new one: say the Malayalam for a discovery, then the Malayalam for scientific.
\end{quote}

\subsection*{You'll want to know: \ml{ഗണിതം}}
\textbf{\ml{ഗണിതം}} — \emph{gaṇitaṁ} — \textquotedblleft{}formal mathematics; maths\textquotedblright{}.

\begin{quote}
\ml{എനിക്ക്} \ml{ഗണിതം} \ml{ഒട്ടും} \ml{ഇഷ്ടമല്ല}.
\end{quote}

\begin{quote}
\emph{enikkŭ gaṇitaṁ oṭṭuṁ iṣṭamalla.}
\end{quote}

\begin{quote}
I don't like mathematics at all.
\end{quote}

At school people also say \ml{കണക്ക്}.

\begin{cousinweb}
\textbf{\ml{ഗണിതം}} (\emph{gaṇitaṁ}) is borrowed from Sanskrit \emph{gaṇita}, \textquotedblleft{}counted, reckoned\textquotedblright{} (CDIAL 3996): mathematics is reckoning.
\end{cousinweb}

\begin{grammarlens}[title={What you've built}]
One more word for giving reasons and opinions.
\end{grammarlens}

\subsection*{Your turn}
\begin{itemize}
\raggedright
  \item \emph{Say it:} \emph{gaṇitaṁ}
  \item \emph{Say it:} \emph{gaṇitaṁ}, once more
  \item \emph{Say it:} \emph{enikkŭ gaṇitaṁ oṭṭuṁ iṣṭamalla.} — I don't like mathematics at all.
  \item \emph{Recall:} say the Malayalam for a discovery, then the Malayalam for scientific, then say \emph{gaṇitaṁ} again
\end{itemize}

\begin{tcolorbox}[breakable,colback=teal!4,colframe=teal!35!black,title={Before you move on}]
What does \ml{ഗണിതം} mean? (\textquotedblleft{}Formal mathematics; maths\textquotedblright{}.) Say it once more.
\end{tcolorbox}
\section[\texorpdfstring{\ml{സാമ്പത്തികശാസ്ത്രം} (sāmpattikaśāstraṁ)}{സാമ്പത്തികശാസ്ത്രം (sāmpattikaśāstraṁ)}]{\ml{സാമ്പത്തികശാസ്ത്രം} (sāmpattikaśāstraṁ) — economics}

\label{lesson:ML-C437-sampattikasastram}



\begin{quote}
Before the new one: say the Malayalam for scientific, then the Malayalam for formal mathematics.
\end{quote}

\subsection*{You'll want to know: \ml{സാമ്പത്തികശാസ്ത്രം}}
\textbf{\ml{സാമ്പത്തികശാസ്ത്രം}} — \emph{sāmpattikaśāstraṁ} — \textquotedblleft{}economics\textquotedblright{}.

\begin{quote}
\ml{അവൾ} \ml{കോളേജിൽ} \ml{സാമ്പത്തികശാസ്ത്രം} \ml{പഠിക്കുന്നു}.
\end{quote}

\begin{quote}
\emph{avaḷ kōḷējil sāmpattikaśāstraṁ paṭhikkunnu.}
\end{quote}

\begin{quote}
She studies economics at college.
\end{quote}

\ml{സാമ്പത്തിക}, financial, and \ml{ശാസ്ത്രം}, a science.

\begin{cousinweb}
\textbf{\ml{സാമ്പത്തികശാസ്ത്രം}} (\emph{sāmpattikaśāstraṁ}) is a Malayalam compound: \ml{സാമ്പത്തിക}, \textquotedblleft{}financial\textquotedblright{}, built on Sanskrit \emph{sampatti}, \textquotedblleft{}prosperity, wealth\textquotedblright{} (Monier-Williams), and \ml{ശാസ്ത്രം}, from Sanskrit \emph{śāstra}, \textquotedblleft{}a treatise, a science\textquotedblright{} (CDIAL 12422).
\end{cousinweb}

\begin{grammarlens}[title={What you've built}]
One more word for giving reasons and opinions.
\end{grammarlens}

\subsection*{Your turn}
\begin{itemize}
\raggedright
  \item \emph{Say it:} \emph{sāmpattikaśāstraṁ}
  \item \emph{Say it:} \emph{sāmpattikaśāstraṁ}, once more
  \item \emph{Say it:} \emph{avaḷ kōḷējil sāmpattikaśāstraṁ paṭhikkunnu.} — She studies economics at college.
  \item \emph{Recall:} say the Malayalam for scientific, then the Malayalam for formal mathematics, then say \emph{sāmpattikaśāstraṁ} again
\end{itemize}

\begin{tcolorbox}[breakable,colback=teal!4,colframe=teal!35!black,title={Before you move on}]
What does \ml{സാമ്പത്തികശാസ്ത്രം} mean? (\textquotedblleft{}Economics\textquotedblright{}.) Say it once more.
\end{tcolorbox}
\section[\texorpdfstring{\ml{തത്വശാസ്ത്രം} (tatvaśāstraṁ)}{തത്വശാസ്ത്രം (tatvaśāstraṁ)}]{\ml{തത്വശാസ്ത്രം} (tatvaśāstraṁ) — philosophy}

\label{lesson:ML-C437-tatvasastram}



\begin{quote}
Before the new one: say the Malayalam for formal mathematics, then the Malayalam for economics.
\end{quote}

\subsection*{You'll want to know: \ml{തത്വശാസ്ത്രം}}
\textbf{\ml{തത്വശാസ്ത്രം}} — \emph{tatvaśāstraṁ} — \textquotedblleft{}philosophy\textquotedblright{}.

\begin{quote}
\ml{പ്രൊഫസർ} \ml{തത്വശാസ്ത്രത്തെക്കുറിച്ച്} \ml{ഒരു} \ml{പുസ്തകം} \ml{എഴുതി}.
\end{quote}

\begin{quote}
\emph{prophasar tatvaśāstrattekkuṟiccŭ oru pustakaṁ eḻuti.}
\end{quote}

\begin{quote}
The professor wrote a book on philosophy.
\end{quote}

A philosopher is \ml{തത്വചിന്തകൻ}, one who thinks about principles.

\begin{cousinweb}
\textbf{\ml{തത്വശാസ്ത്രം}} (\emph{tatvaśāstraṁ}) is a Malayalam compound built from Sanskrit \emph{tattva}, \textquotedblleft{}the real state, truth\textquotedblright{} (CDIAL 5642), and \emph{śāstra}, \textquotedblleft{}a treatise, a science\textquotedblright{} (CDIAL 12422): the science of what is real.
\end{cousinweb}

\begin{grammarlens}[title={What you've built}]
One more word for giving reasons and opinions.
\end{grammarlens}

\subsection*{Your turn}
\begin{itemize}
\raggedright
  \item \emph{Say it:} \emph{tatvaśāstraṁ}
  \item \emph{Say it:} \emph{tatvaśāstraṁ}, once more
  \item \emph{Say it:} \emph{prophasar tatvaśāstrattekkuṟiccŭ oru pustakaṁ eḻuti.} — The professor wrote a book on philosophy.
  \item \emph{Recall:} say the Malayalam for formal mathematics, then the Malayalam for economics, then say \emph{tatvaśāstraṁ} again
\end{itemize}

\begin{tcolorbox}[breakable,colback=teal!4,colframe=teal!35!black,title={Before you move on}]
What does \ml{തത്വശാസ്ത്രം} mean? (\textquotedblleft{}Philosophy\textquotedblright{}.) Say it once more.
\end{tcolorbox}
\section[\texorpdfstring{\ml{വിജ്ഞാനം} (vijñānaṁ)}{വിജ്ഞാനം (vijñānaṁ)}]{\ml{വിജ്ഞാനം} (vijñānaṁ) — knowledge; learning}

\label{lesson:ML-C437-vijnanam}



\begin{quote}
Before the new one: say the Malayalam for economics, then the Malayalam for philosophy.
\end{quote}

\subsection*{You'll want to know: \ml{വിജ്ഞാനം}}
\textbf{\ml{വിജ്ഞാനം}} — \emph{vijñānaṁ} — \textquotedblleft{}knowledge; learning\textquotedblright{}.

\begin{quote}
\ml{പുസ്തകങ്ങൾ} \ml{വിജ്ഞാനത്തിന്റെ} \ml{ഉറവിടമാണ്}.
\end{quote}

\begin{quote}
\emph{pustakaṅṅaḷ vijñānattinṟe uṟaviṭamāṇŭ.}
\end{quote}

\begin{quote}
Books are a source of knowledge.
\end{quote}

An encyclopedia is \ml{വിജ്ഞാനകോശം}, a treasury of knowledge.

\begin{cousinweb}
\textbf{\ml{വിജ്ഞാനം}} (\emph{vijñānaṁ}) is borrowed from Sanskrit \emph{vijñāna}, \textquotedblleft{}the act of discerning, understanding\textquotedblright{} (CDIAL 11707), from the root \emph{jñā}, \textquotedblleft{}to know\textquotedblright{}.
\end{cousinweb}

\begin{grammarlens}[title={What you've built}]
One more word for giving reasons and opinions.
\end{grammarlens}

\subsection*{Your turn}
\begin{itemize}
\raggedright
  \item \emph{Say it:} \emph{vijñānaṁ}
  \item \emph{Say it:} \emph{vijñānaṁ}, once more
  \item \emph{Say it:} \emph{pustakaṅṅaḷ vijñānattinṟe uṟaviṭamāṇŭ.} — Books are a source of knowledge.
  \item \emph{Recall:} say the Malayalam for economics, then the Malayalam for philosophy, then say \emph{vijñānaṁ} again
\end{itemize}

\begin{tcolorbox}[breakable,colback=teal!4,colframe=teal!35!black,title={Before you move on}]
What does \ml{വിജ്ഞാനം} mean? (\textquotedblleft{}Knowledge; learning\textquotedblright{}.) Say it once more.
\end{tcolorbox}
\section[\texorpdfstring{\ml{പഠനം} (paṭhanaṁ)}{പഠനം (paṭhanaṁ)}]{\ml{പഠനം} (paṭhanaṁ) — studies; a study}

\label{lesson:ML-C437-pathanam}



\begin{quote}
Before the new one: say the Malayalam for philosophy, then the Malayalam for knowledge.
\end{quote}

\subsection*{You'll want to know: \ml{പഠനം}}
\textbf{\ml{പഠനം}} — \emph{paṭhanaṁ} — \textquotedblleft{}studies; a study\textquotedblright{}.

\begin{quote}
\ml{അവന്റെ} \ml{പഠനം} \ml{കഴിഞ്ഞ} \ml{വർഷം} \ml{തീർന്നു}.
\end{quote}

\begin{quote}
\emph{avanṟe paṭhanaṁ kaḻiñña varṣaṁ tīrnnu.}
\end{quote}

\begin{quote}
His studies finished last year.
\end{quote}

The noun of \ml{പഠിക്കുക}, to learn. A study, a piece of research, is also a \ml{പഠനം}.

\begin{cousinweb}
\textbf{\ml{പഠനം}} (\emph{paṭhanaṁ}) is borrowed from Sanskrit \emph{paṭhana}, \textquotedblleft{}reading\textquotedblright{} (CDIAL 7713), from the root \emph{paṭh}, \textquotedblleft{}to read aloud\textquotedblright{}, the root of \ml{പഠിക്കുക}.
\end{cousinweb}

\begin{grammarlens}[title={What you've built}]
One more word for giving reasons and opinions.
\end{grammarlens}

\subsection*{Your turn}
\begin{itemize}
\raggedright
  \item \emph{Say it:} \emph{paṭhanaṁ}
  \item \emph{Say it:} \emph{paṭhanaṁ}, once more
  \item \emph{Say it:} \emph{avanṟe paṭhanaṁ kaḻiñña varṣaṁ tīrnnu.} — His studies finished last year.
  \item \emph{Recall:} say the Malayalam for philosophy, then the Malayalam for knowledge, then say \emph{paṭhanaṁ} again
\end{itemize}

\begin{tcolorbox}[breakable,colback=teal!4,colframe=teal!35!black,title={Before you move on}]
What does \ml{പഠനം} mean? (\textquotedblleft{}Studies; a study\textquotedblright{}.) Say it once more.
\end{tcolorbox}
